\subsection{3. Thảo luận chung về kết quả thực nghiệm}

\indent\indent Dựa trên các thực nghiệm đối sánh từ mô hình mạng nơ-ron đồ thị cơ sở (GNN baseline) đến kiến trúc phân lớp tách rời (\textit{Decoupled GNN-ML}), một số phát hiện và thảo luận tổng quát về đặc tính của hệ thống được đúc kết như sau:

\begin{itemize}[label={--}, itemsep=1pt]
    \item \textbf{Vai trò của đồ thị đa tầng nhận thức ràng buộc:} Sự thành công của các thuật toán không chỉ đến từ kiến trúc mạng mà còn phụ thuộc quyết định vào chất lượng không gian biểu diễn ban đầu. Việc nén 18 loại quan hệ phức tạp thành 4 tầng đồ thị (Follow, Interact, Similarity, Co-Owner), kết hợp cùng cơ chế chuẩn hóa trọng số logarit, đã giải quyết triệt để bài toán bùng nổ thông điệp. Lớp ràng buộc nhận thức ($w_{ij}$) đóng vai trò như một bộ lọc định lượng, ép các mô hình GNN phải đánh giá láng giềng dựa trên chất lượng hành vi thay vì bị phân tâm vào số lượng quan hệ giữa 2 nút.

    \item \textbf{Hiệu quả của phương pháp gán nhãn giả định dựa trên hệ thống luật:} Trong bối cảnh thiếu hụt nghiêm trọng các nhãn thực tế, hệ thống đã chứng minh tính đúng đắn của quy trình trích xuất đặc trưng ẩn (GAE) và phân cụm (K-Means). Quy trình kết hợp với cơ chế kiểm chứng dữ liệu nhãn không chỉ giải quyết bài toán về dữ liệu mà còn mở ra một hướng tiếp cận để tự động hóa việc dán nhãn trên dữ liệu của các mạng xã hội phi tập trung quy mô lớn.

    \item \textbf{Hiệu quả biểu diễn đặc trưng của cơ chế Attention:} Dựa trên nền tảng đồ thị đa tầng, kết quả thực nghiệm khẳng định các kiến trúc tích hợp cơ chế chú ý (tiêu biểu là GATv2) duy trì được hiệu năng ổn định hơn đáng kể so với phương pháp tổng hợp tĩnh (GCN) khi đồ thị mở rộng.
    
    \item \textbf{Khả năng kháng nhiễu và tính linh hoạt của Decoupled GNN-ML:} Việc phân tách quá trình biểu diễn đồ thị và phân lớp bộc lộ ưu thế rõ rệt trên tập dữ liệu lớn (phạm vi Quý). Khi các mô hình GNN cơ sở có dấu hiệu bão hòa do mật độ liên kết quá dày đặc, tổ hợp GATv2 + Random Forest vẫn thiết lập được ranh giới phi tuyến tính vững chắc.

    \item \textbf{Tính ứng dụng và sự an toàn trong môi trường thực tiễn:} Bên cạnh chỉ số Macro F1-Score vượt trội (0.9303), giá trị cốt lõi của mô hình nằm ở khả năng kiểm soát tỷ lệ False Positive. Việc hệ thống ghi nhận tỷ lệ 0\% nhận diện nhầm từ người dùng thật sang tài khoản độc hại trên tập thử nghiệm Quý là một bảo chứng về độ an toàn. Tính chất này giải quyết hoàn hảo bài toán đánh đổi giữa việc thanh lọc mạng lưới và bảo vệ tuyệt đối trải nghiệm của người dùng chân chính.
    
    \item \textbf{Một số giới hạn và thách thức còn tồn tại:} Mặc dù đạt hiệu năng cao, mô hình vẫn ghi nhận sự giao thoa dự đoán giữa nhóm \texttt{HIGH\_RISK} và \texttt{MALICIOUS}. Điều này xuất phát từ bản chất động của các mạng lưới có quy mô lớn, nơi ranh giới hành vi chưa hoàn toàn tách biệt tại thời điểm cắt lớp dữ liệu. Ngoài ra, do giới hạn của phương pháp giám sát yếu, tập nhãn giả định dù đã qua hiệu chuẩn vẫn có thể tiềm ẩn một tỷ lệ nhiễu cấu trúc nhỏ.
\end{itemize}